\section{余法丛}
\label{sec:conormal-bundles}
\stfhIndex{I185}{la^ge^lang^ri^zi^liu^xing}{拉格朗日子流形（Lagrangian submanifold）}
\stfhIndex{I209}{zi^ran^bian^jie^tiao^jian}{自然边界条件（natural boundary condition）}

在余切丛 $T^*N$ 上，哈密顿作用量泛函定义为积分
\[
  \mathcal A_H(\gamma)=\int_\gamma\theta-\int_0^1H(t,\gamma(t))\,dt.
\]
通常把 $\theta$ 写作 $p\,dq$；我们也会不时采用这一写法。下面推导 $\mathcal A_H$ 的第一变分公式。沿给定道路 $\gamma$ 取变分 $\xi$，其支集包含在 $[0,1]$ 的某个小开子集中，并且相应的像包含在 $T^*U$ 中，其中 $U\subset N$ 是一个坐标邻域；这样便可使用典范坐标来计算。此时
$\theta=\sum p_j\,dq^j$，因而 $\mathcal A_H$ 与经典相空间上的哈密顿作用量泛函（1.3.26）具有完全相同的形式。更形式化地，（1.3.26）写成
\begin{equation}
\begin{split}
 d\mathcal A_H(\gamma)(\xi)
  ={}&\langle\theta(\gamma(1)),\xi(1)\rangle
      -\langle\theta(\gamma(0)),\xi(0)\rangle \\
    &+\int_0^1\omega_0\bigl(\dot\gamma-X_H(\gamma),\xi\bigr)\,dt.
\end{split}
\label{eq:conormal-first-variation}
\end{equation}

\begin{exercise}\label{ex:conormal-first-variation}
证明这个公式。
\end{exercise}

这里回顾一下：在我们的约定中，余切丛上的典范辛形式 $\omega_0$ 由 $\omega_0=-d\theta$ 给出。

在经典相空间中，除了周期边界条件以外，我们还考察过作为自然边界条件的两点边界条件；也就是说，它是使\eqref{eq:conormal-first-variation}中的边界项消失的边界条件。对给定的两个位形点 $x,y\in N$，这相当于把 $\mathcal A_H$ 限制在道路空间
\[
 \mathcal P(T^*N;T_x^*N,T_y^*N)
 =\{\gamma:[0,1]\to T^*N\mid
       \gamma(0)\in T_x^*N,\ \gamma(1)\in T_y^*N\}
\]
上。事实表明，这种自然边界条件容许一个比两点边界条件更一般的类别，我们称之为\emph{余法边界条件}。
\stfhIndex{I083}{yu^fa^bian^jie^tiao^jian}{余法边界条件（conormal boundary condition）}

\begin{definition}[余法丛]\label{def:conormal-bundle}
\stfhIndex{I084}{yu^fa^cong}{余法丛（conormal bundle）}
设 $S\subset N$ 是光滑子流形。$S$ 的\emph{余法丛}定义为
\[
 \nu^*S=\{\alpha\in T_q^*N\mid q\in S,\ \alpha|_{T_qS}=0\}.
\]
\end{definition}

下面的引理给出余法丛的基本性质。它说明余法边界条件确是作用量泛函 $\mathcal A_H$ 的\emph{自然边界条件}，因为这一条件使\eqref{eq:conormal-first-variation}中的边界贡献消失。

\begin{lemma}\label{lem:conormal-tautological-vanishing}
设 $S\subset N$ 是光滑子流形，$i_S:\nu^*S\hookrightarrow T^*N$ 是其余法丛的包含映射。则
$i_S^*\theta=0$。特别地，$i_S^*\omega_0=0$。此外，对任意 $S\subset N$，均有
$\dim\nu^*S=n\;(=\dim N)$。
\end{lemma}

\begin{proof}
设 $\xi\in T_\alpha(\nu^*S)$，则 $d\pi(\xi)\in T_{\pi(\alpha)}S$，所以
\[
 \theta_\alpha(\xi)=\alpha(d\pi(\xi))=0,
\]
因为 $\alpha\in\nu^*_{\pi(\alpha)}S$。
\end{proof}

余法边界条件包含如下两个特例。

\begin{example}\label{ex:conormal-special-cases}
\begin{enumerate}[label=(\arabic*)]
\item 设 $S=\{q\}$ 是 $N$ 中的一点，则 $\nu^*\{q\}=T_q^*N$，即 $T^*N$ 在 $q$ 上的纤维。
\item 当 $S=N$ 时，$\nu^*S=o_N$，即零截面。
\end{enumerate}
\end{example}

现在设 $S_0,S_1\subset N$ 是两个子流形，$\nu^*S_i$（$i=0,1$）是它们的余法丛。考虑 $T^*N$ 上道路组成的子集
\[
 \mathcal L(\nu^*S_0,\nu^*S_1)
 =\{\gamma:[0,1]\to T^*N\mid
       \gamma(0)\in\nu^*S_0,\ \gamma(1)\in\nu^*S_1\}.
\]
用 $d\mathcal A_{(H;S_0,S_1)}$ 表示 $d\mathcal A_H$ 在
$\mathcal L(\nu^*S_0,\nu^*S_1)$ 上的限制，则
\[
 d\mathcal A_{(H;S_0,S_1)}(\gamma)(\xi)
 =\int_0^1\omega\bigl(\dot\gamma-X_H(t,\gamma(t)),\xi(t)\bigr)\,dt.
\]
因此，临界点集 $\operatorname{Crit}(\mathcal A_{(H;S_0,S_1)})$ 与从
$\nu^*S_0$ 到 $\nu^*S_1$ 的哈密顿轨道一一对应；换言之，它与满足
\[
 \dot z=X_H(t,z),\qquad z(0)\in\nu^*S_0,\quad z(1)\in\nu^*S_1
\]
的道路 $z$ 一一对应。

在一般辛流形 $(M,\omega)$ 上，并没有作用量泛函 $\mathcal A_H$ 的直接类似物，因为它的定义本质上使用了 $1$-形式 $\theta$。为了把上述哈密顿方程看成 $M$ 的道路空间上某个泛函的极值道路，我们首先需要发展一般辛流形的大量背景几何。我们从若干基本的辛线性代数开始。
